\chapter{Il modulo \code{runtime-config}}
\label{cap:14-runtime-config}

Il modulo \code{runtime-config} consente di modificare a caldo un sottoinsieme
dei parametri della piattaforma, senza riavviare il processo. Sono 540 righe in
dieci classi.

Il modulo ha una motivazione strettamente sperimentale. Un esperimento che
confronta due configurazioni deve, per essere valido, mantenere costante tutto il
resto: stesso processo, stessa JVM riscaldata, stesse code popolate, stesso stato
del garbage collector. Riavviare il control plane fra una configurazione e
l'altra reintroduce esattamente le variabili che l'esperimento intende
escludere.

\section{Che cosa \`e modificabile a caldo}

Sei parametri, esattamente:

\begin{lstlisting}[language=Java]
public record RuntimeConfigSnapshot(
        long revision,
        int rateMaxPerSecond,
        boolean syncQueueEnabled,
        boolean syncQueueAdmissionEnabled,
        Duration syncQueueMaxEstimatedWait,
        Duration syncQueueMaxQueueWait,
        int syncQueueRetryAfterSeconds
) { ... }
\end{lstlisting}

Il criterio di inclusione \`e la reversibilit\`a senza effetti sullo stato
esistente. Il limite di frequenza \`e un intero letto a ogni richiesta:
cambiarlo ha effetto dalla richiesta successiva e non invalida nulla. Le soglie
della coda sincrona sono lette a ogni valutazione di ammissione (\cref{cap:10-sync-queue}):
lo stesso.

Sono invece esclusi \code{max-depth} della coda sincrona --- ridurla
richiederebbe di decidere che cosa fare degli elementi in eccesso --- e la
selezione dei moduli, che per costruzione avviene a tempo di build
(\cref{cap:08-sistema-moduli}).

\section{Il meccanismo: snapshot immutabile e revisione}

Lo stato \`e un unico record immutabile in un \code{AtomicReference}. Ogni
aggiornamento ne produce uno nuovo con revisione incrementata. I consumatori
leggono sempre l'ultimo:

\begin{lstlisting}[language=Java]
@Override
public Duration syncQueueMaxEstimatedWait() {
    return current.get().syncQueueMaxEstimatedWait();
}
\end{lstlisting}

La propriet\`a che questa costruzione garantisce \`e che una singola valutazione
di ammissione veda una configurazione \emph{coerente}: legge uno snapshot, e i
campi di quello snapshot appartengono tutti alla stessa revisione. Con sei
variabili volatili indipendenti, un aggiornamento concorrente potrebbe far s\`i
che una valutazione legga la nuova soglia di attesa e il vecchio flag di
ammissione --- una combinazione che nessuno ha mai richiesto.

\subsection{Il lock ottimistico}

L'aggiornamento richiede al chiamante di dichiarare la revisione che si aspetta:

\begin{lstlisting}[language=Java]
public RuntimeConfigSnapshot update(long expectedRevision, RuntimeConfigPatch patch) {
    while (true) {
        RuntimeConfigSnapshot snapshot = current.get();
        if (snapshot.revision() != expectedRevision) {
            throw new RevisionMismatchException(expectedRevision, snapshot.revision());
        }
        RuntimeConfigSnapshot candidate = snapshot.applyPatch(patch);
        if (current.compareAndSet(snapshot, candidate)) {
            return candidate;
        }
        // CAS failed because another thread updated; retry loop will re-check revision
    }
}
\end{lstlisting}

\`E il pattern \emph{compare-and-swap} classico, con il controllo di revisione
che lo trasforma in un lock ottimistico visibile al chiamante. Un aggiornamento
basato su uno stato obsoleto riceve \code{409 Conflict} con la revisione
corrente, anzich\'e sovrascrivere in silenzio una modifica concorrente.

L'aggiornamento \`e un \code{PATCH} con semantica ``i campi nulli restano
invariati''. La combinazione di patch parziale e revisione attesa \`e ci\`o che
rende sicure due modifiche concorrenti su campi diversi: la seconda vede il
conflitto, rilegge, e riapplica.

\section{L'API di amministrazione}

\begin{center}
\small
\begin{adjustbox}{max width=\textwidth}
\begin{tabular}{@{}llp{6.6cm}@{}}
\toprule
\textbf{Metodo} & \textbf{Percorso} & \textbf{Semantica} \\
\midrule
\code{GET} & \code{/v1/admin/runtime-config} & Snapshot corrente con revisione. \\
\code{POST} & \code{/v1/admin/runtime-config/validate} & Prova la patch senza
applicarla. \\
\code{PATCH} & \code{/v1/admin/runtime-config} & Applica, con
\code{expectedRevision} obbligatoria. \\
\bottomrule
\end{tabular}
\end{adjustbox}
\end{center}

\subsection{L'endpoint di validazione}

L'esistenza di \code{/validate} merita una nota. Applica la patch a una copia
dello snapshot corrente, ne verifica la validit\`a e riporta gli errori, senza
modificare nulla. \`E lo stesso codice che il \code{PATCH} esegue prima di
applicare.

Il suo valore \`e per gli strumenti automatici: uno script di esperimento pu\`o
verificare che la propria matrice di configurazioni sia interamente valida prima
di iniziare, anzich\'e scoprire al terzo punto della matrice che una combinazione
verr\`a rifiutata --- avendo gi\`a consumato il tempo dei primi due.

\subsection{La sequenza del \code{PATCH}}

L'ordine delle operazioni codifica una scelta di sicurezza:

\begin{enumerate}[leftmargin=*]
  \item \code{expectedRevision} presente, altrimenti \code{400};
  \item la patch \`e sintatticamente valida (durate parsabili, ecc.), altrimenti
  \code{400} con nome del campo e valore offensivo;
  \item lo snapshot risultante \`e semanticamente valido, altrimenti \code{422};
  \item la revisione corrisponde, altrimenti \code{409} con la revisione reale;
  \item lo snapshot viene sostituito atomicamente;
  \item l'applicatore propaga ai componenti che hanno bisogno di essere
  informati, con rollback in caso di fallimento (\code{503}).
\end{enumerate}

La validazione precede la sostituzione: uno snapshot invalido non entra mai nello
stato. Il rollback al punto 6 copre il caso residuo in cui la propagazione
fallisca dopo che lo stato \`e cambiato.

\subsection{Il metodo sincronizzato}

\begin{lstlisting}[language=Java]
@PatchMapping
public synchronized ResponseEntity<Object> patch(@RequestBody PatchRequest request) { ... }
\end{lstlisting}

Il metodo del controller \`e \code{synchronized}, il che serializza tutte le
modifiche di configurazione. \`E ridondante rispetto al CAS del servizio per la
correttezza dello stato, ma serializza anche la fase di applicazione, che il CAS
non copre: due patch concorrenti potrebbero altrimenti propagare i propri effetti
in ordine invertito rispetto a quello in cui hanno vinto il CAS.

\`E anche una violazione del principio non bloccante del
\cref{cap:03-requisiti-principi}, ed \`e accettabile perch\'e riguarda un
endpoint amministrativo a bassissima frequenza. Vale la pena notare che il
codice non lo giustifica, e che il lettore deve dedurre da s\'e che il costo \`e
confinato.

\section{L'applicatore}

\code{RuntimeConfigApplier} propaga lo snapshot ai componenti mutabili in ordine
deterministico:

\begin{lstlisting}[language=Java]
/**
 * Applies the new snapshot to all mutable runtime components.
 * Sync queue fields are read directly from {@link RuntimeConfigService} by consumers,
 * so only components with dedicated setters need explicit application here.
 */
public void apply(RuntimeConfigSnapshot snapshot, RuntimeConfigSnapshot previous,
                  RuntimeConfigService configService) { ... }
\end{lstlisting}

Il commento rivela che esistono due modalit\`a di propagazione:

\begin{description}[style=nextline, leftmargin=0pt]
  \item[\textbf{Pull}] --- i consumatori della coda sincrona leggono dallo
  snapshot a ogni uso. Non c'\`e nulla da propagare; la modifica ha effetto
  immediatamente e senza coordinamento.
  \item[\textbf{Push}] --- il limitatore di frequenza possiede una propria
  variabile interna, e l'applicatore deve chiamarne il setter.
\end{description}

La modalit\`a pull \`e evidentemente preferibile e il modulo la usa dove
possibile; la push esiste per un componente del nucleo che \`e antecedente al
modulo e che il modulo non deve costringere a cambiare struttura. \`E una
concessione ragionevole al principio di non far dipendere il nucleo dai moduli.

\section{Osservabilit\`a}

\begin{center}
\small
\begin{adjustbox}{max width=\textwidth}
\begin{tabular}{@{}lll@{}}
\toprule
\textbf{Metrica} & \textbf{Tipo} & \textbf{Significato} \\
\midrule
\code{controlplane\_runtime\_config\_revision} & gauge & Revisione corrente. \\
\code{controlplane\_runtime\_config\_updates\_total} & counter &
Aggiornamenti, con etichetta \code{status}. \\
\code{controlplane\_runtime\_config\_apply\_duration\_seconds} & timer &
Durata dell'applicazione. \\
\bottomrule
\end{tabular}
\end{adjustbox}
\end{center}

Il gauge sulla revisione \`e lo strumento pi\`u utile del gruppo in contesto
sperimentale: esportando la revisione nelle stesse serie temporali delle metriche
di prestazione, un cambio di configurazione diventa un gradino osservabile sul
grafico, e l'attribuzione di una variazione di latenza al cambio che l'ha causata
non richiede di correlare a mano i log con le misure.

\section{Disattivato per default}

\begin{lstlisting}[language=yamlnf]
nanofaas:
  admin:
    runtime-config:
      enabled: false
\end{lstlisting}

L'API espone operazioni di scrittura su una piattaforma priva di autenticazione
(\cref{cap:03-requisiti-principi}). Il default disattivato fa s\`i che
l'esposizione sia una decisione consapevole dell'operatore anzich\'e una
propriet\`a ereditata dalla scelta dei moduli. \`E l'unico modulo che, pur essendo
compilato nel binario, resta inerte finch\'e non viene attivato esplicitamente.
